% Source: Wald GR, physical PDF pp. 202--208
%         (printed pp. 189--195).
% Coverage: Section 8.1, Futures and Pasts: Basic Definitions and Results;
%           equations (8.1.1)--(8.1.3), lemmas 8.1.1, 8.1.4--8.1.5,
%           theorems 8.1.2--8.1.3 and 8.1.6, and the unnumbered corollary.
% Figures/tables/footnotes: Figures 8.1--8.5; footnotes 1--6.
% Status: source-reviewed and compile-clean.
% Uncertainties: none.

\subsection{Futures and Pasts: Basic Definitions and Results}\label{sec:8.1}
\setcounter{theorem}{0}

Let \((M,g_{ab})\) be a spacetime. At each event \(p\in M\), the tangent space,
\(T_pM\), is, of course, isomorphic to Minkowski spacetime. We will refer to the light
cone passing through the origin of \(T_pM\) as the \emph{light cone of \(p\)}. Thus, we
emphasize that the light cone of \(p\) is a subset of \(T_pM\), not \(M\).\footnote{It
should be noted that some authors use the term \lq\lq light cone\rq\rq\ to designate the subset
of \(M\) generated by null geodesics from \(p\) (as, in fact, we did in chapter 1). They
use the term \lq\lq null cone\rq\rq\ to designate what we have called the \lq\lq light cone.\rq\rq} As in
special relativity, at each \(p\in M\) we may designate half of the light cone as
\lq\lq future\rq\rq\ and the other half as \lq\lq past.\rq\rq\ However, in a non-simply
connected\footnote{See chapter 13 for the definition of a non-simply connected manifold.}
spacetime it may not be possible to make a continuous designation of \lq\lq future\rq\rq\ and
\lq\lq past\rq\rq\ as \(p\) varies over \(M\). An example of a spacetime for which no such
continuous designation can be made is shown in Fig.~\ref{fig:8.1}. If a continuous
choice can be made, \((M,g_{ab})\) is said to be \emph{time orientable}. (This property
of a spacetime is analogous to, but distinct from, the notion of orientability of a
manifold defined in section B.2 of appendix B.) Thus, non-time orientable spacetimes
have the physically pathological property that we cannot consistently distinguish
between the notions of going \lq\lq forward in time\rq\rq\ as opposed to \lq\lq backward in time.\rq\rq\
In the following, we will consider only time orientable spacetimes and will assume
that a continuous designation has been made of the \lq\lq future\rq\rq\ and \lq\lq past\rq\rq\ halves of
the light cones at each point. A timelike or null vector lying in the \lq\lq future half\rq\rq\
of the light cone will be called \emph{future directed}.

% Source: Wald GR, physical PDF p. 202
%         (printed p. 189).
% Coverage: A non-time-orientable spacetime.
% Figures/tables/footnotes: Figure 8.1; no tables or footnotes.
% Status: source-reviewed and compile-clean.
% Uncertainties: none.

\begingroup
\renewcommand{\thefigure}{8.1}
\begin{figure}[H]
  \centering
  \begin{tikzpicture}[x=.92cm,y=.92cm,>=Latex,line cap=round,font=\small]
    \draw (0,0) -- (0,4.05);
    \draw (6,0) -- (6,3.95);
    \draw[->,very thick] (0,1.65) -- (0,2.45);
    \draw[->,very thick] (6,1.65) -- (6,2.45);
    \draw[<->] (.12,3.55) -- (5.88,3.55) node[midway,above] {identify};
    \foreach \x/\a in {.68/0,1.82/56,3.00/90,4.18/-56,5.32/0} {
      \begin{scope}[shift={(\x,1.64)},rotate=\a]
        \draw (-.22,-.38) -- (0,0) -- (-.22,.38);
        \draw (.22,-.38) -- (0,0) -- (.22,.38);
        \draw (-.22,-.38) .. controls (-.04,-.42) and (.04,-.42) .. (.22,-.38);
        \draw (-.22,.38) .. controls (-.04,.42) and (.04,.42) .. (.22,.38);
      \end{scope}
    }
  \end{tikzpicture}
  \caption{A non-time-orientable spacetime.}
  \label{fig:8.1}
\end{figure}
\endgroup


An important property satisfied by every time orientable spacetime is expressed in
the following lemma:

\begin{lemma}\label{lem:8.1.1}
Let \((M,g_{ab})\) be time orientable. Then there exists a (highly nonunique) smooth
nonvanishing timelike vector field \(t^a\) on \(M\).
\end{lemma}
\begin{proof}
Since \(M\) is paracompact, we can choose a smooth Riemannian metric \(k_{ab}\) on
\(M\) (see appendix A). At each \(p\in M\) there will be a unique future directed
timelike vector \(t^a\) which minimizes the value of \(g_{ab}v^av^b\) for vectors
\(v^a\) subject to the condition that \(k_{ab}v^av^b=1\). This \(t^a\) will vary
smoothly over \(M\) and thus provide the desired vector field.
\end{proof}

Conversely, if a continuous, timelike vector field can be chosen, then \((M,g_{ab})\)
is time orientable.

Let \((M,g_{ab})\) be a time orientable spacetime. A differentiable curve
\(\lambda(t)\) is said to be a \emph{future directed timelike curve} if at each
\(p\in\lambda\) the tangent \(t^a\) is a future directed timelike vector. Note that
according to this definition, if \(t^a\) vanishes at one point, \(\lambda\) is not
considered to be a timelike curve.\footnote{The strictly timelike character of
\(t^a\) is imposed in order that the curve \(\lambda(t)=p\) for all \(t\) is not
considered to be timelike. Note that this makes the notion of a differentiable
timelike curve parameterization dependent. However, below we shall define a
parameterization-independent notion of a continuous, future-directed timelike curve.}
Similarly, \(\lambda\) is said to be a \emph{future directed causal curve} if at each
\(p\in\lambda\), \(t^a\) is either a future directed timelike or null vector. Thus,
on a causal curve \(t^a\) may vanish. Analogous definitions apply to past directed
timelike and causal curves.

The \emph{chronological future} of \(p\in M\), denoted \(I^+(p)\), is defined as the
set of events that can be reached by a future directed timelike curve starting from
\(p\),
\begin{equation}
  I^+(p)=\left\{q\in M\ \middle|\ \begin{gathered}
    \text{There exists a future directed timelike curve }\lambda(t)\\
    \text{with }\lambda(0)=p\text{ and }\lambda(1)=q
  \end{gathered}\right\}.
  \tag{8.1.1}\label{eq:8.1.1}
\end{equation}
Since one always can perform a sufficiently small deformation of the endpoint of a
timelike curve while preserving the timelike nature of the curve, it follows that for
all \(q\in I^+(p)\) there exists an open neighborhood \(O\) of \(q\) such that
\(O\subseteq I^+(p)\). (A formal proof of this statement can be given using
theorem~\ref{thm:8.1.2} below.) Thus, \(I^+(p)\) always is an open subset of \(M\).
Note that in general \(p\notin I^+(p)\), but \(p\) will be in \(I^+(p)\) if there is
a closed timelike curve beginning and ending at \(p\).

For any subset \(S\subseteq M\), we define \(I^+(S)\) by
\begin{equation}
  I^+(S)=\bigcup_{p\in S}I^+(p).
  \tag{8.1.2}\label{eq:8.1.2}
\end{equation}
Since an arbitrary union of open sets is open, it follows that \(I^+(S)\) always is
an open set. We also have \(I^+(I^+(S))=I^+(S)\) and
\(I^+(\overline S)=I^+(S)\), where \(\overline S\) denotes the closure of \(S\)
(see problem 1). Analogous definitions and properties apply to the
\emph{chronological pasts} \(I^-(p)\) and \(I^-(S)\).

The \emph{causal future} of \(p\in M\), denoted \(J^+(p)\), is defined in the same
way as \(I^+(p)\) except that \lq\lq future directed causal curve\rq\rq\ replaces \lq\lq future
directed timelike curve\rq\rq\ in Eq.~\eqref{eq:8.1.1}. Thus, we always have
\(p\in J^+(p)\). In flat spacetime, \(J^+(p)\) is a closed set. However, in a
general spacetime, \(J^+(p)\) can fail to be closed; an example of such a spacetime
is given in Fig.~\ref{fig:8.2}. In section 8.3 we shall prove that \(J^+(p)\) must be
closed in any globally hyperbolic spacetime. Again, we define
\begin{equation}
  J^+(S)=\bigcup_{p\in S}J^+(p).
  \tag{8.1.3}\label{eq:8.1.3}
\end{equation}
and define the causal pasts, \(J^-(p)\) and \(J^-(S)\), analogously.\footnote{In
other references the notation \(p\ll q\) and \(p\prec q\) is sometimes used,
respectively, for \(q\in I^+(p)\) and \(q\in J^+(p)\).}

In Minkowski spacetime, \(I^+(p)\) consists precisely of the points that can be
reached by future directed timelike geodesics starting from \(p\). The boundary
\(\partial I^+(p)\), of \(I^+(p)\), is generated by the future directed null geodesics
starting from \(p\). In arbitrary spacetimes neither of these statements is valid in
general, as shown by simple examples obtained by removing points from Minkowski
spacetime such as in Fig.~\ref{fig:8.2}. However, \emph{locally} these properties
always remain valid, as expressed by the following theorem:

% Source: Wald GR, physical PDF p. 204
%         (printed p. 191).
% Coverage: Minkowski spacetime with a removed point on a future light cone.
% Figures/tables/footnotes: Figure 8.2; no tables or footnotes.
% Status: source-reviewed and compile-clean.
% Uncertainties: none.

\begingroup
\renewcommand{\thefigure}{8.2}
\begin{figure}[H]
  \centering
  \begin{tikzpicture}[x=.9cm,y=.9cm,>=Latex,line cap=round,font=\small]
    \draw[thick] (0,0) -- (-2.05,3.22);
    \draw[thick] (0,0) -- (2.05,3.22);
    \draw[thick] (-2.05,3.22) arc[start angle=180,end angle=360,x radius=2.05,y radius=.42];
    \draw[thick] (-2.05,3.22) arc[start angle=180,end angle=0,x radius=2.05,y radius=.42];
    \fill (0,0) circle (1.7pt) node[below] {\(p\)};
    \fill (1.72,2.52) circle (1.7pt) node[right] {\(q\)};
    \draw[fill=white] (1.26,1.77) circle (2.2pt);
    \draw[->] (1.93,1.72) .. controls (2.40,1.72) and (2.33,1.48) .. (1.42,1.70);
    \node[align=left,right] at (2.18,1.46) {(remove\\point)};
  \end{tikzpicture}
  \caption{Minkowski spacetime with a point on the future light cone of \(p\)
  removed. In this spacetime, no causal curve connects \(p\) and \(q\) so
  \(q\notin J^+(p)\) but \(q\in\overline{J^+(p)}\). Thus, \(J^+(p)\) is not
  closed.}
  \label{fig:8.2}
\end{figure}
\endgroup


\begin{theorem}\label{thm:8.1.2}
Let \((M,g_{ab})\) be an arbitrary spacetime, and let \(p\in M\). Then there
exists a convex normal neighborhood of \(p\), i.e., an open set \(U\) with
\(p\in U\) such that for all \(q,r\in U\) there exists a unique geodesic \(\gamma\)
connecting \(q\) and \(r\) and staying entirely within \(U\). Furthermore, for any
such \(U\), \(I^+(p)|_U\) consists of all points reached by future directed timelike
geodesics starting from \(p\) and contained within \(U\), where \(I^+(p)|_U\)
denotes the chronological future of \(p\) in the spacetime \((U,g_{ab})\). In
addition, \(\partial I^+(p)|_U\) is generated by the future directed null geodesics
in \(U\) emanating from \(p\).
\end{theorem}

Although theorem~\ref{thm:8.1.2} may seem obvious intuitively, it is nontrivial to
give a formal proof of it. A proof of the first half is given beginning on page 134 of
Hicks (1965), while the second half is proven as proposition 4.5.1 of Hawking and
Ellis (1973).

If \(q\in J^+(p)\) and \(\lambda\) is a causal curve beginning at \(p\) and ending
at \(q\), we can cover \(\lambda\) by convex normal neighborhoods, and using the
compactness of \(\lambda\) (since it is the continuous image of a closed interval),
we can extract a finite subcover (see appendix A). If \(\lambda\) failed to be a null
geodesic in any such neighborhood, then with the help of theorem~\ref{thm:8.1.2}
we could deform \(\lambda\) into a timelike curve in that neighborhood and then
extend this deformation to the other neighborhoods to obtain a timelike curve from
\(p\) to \(q\). Thus, we obtain the following corollary of theorem~\ref{thm:8.1.2},
which will be strengthened by theorem 9.3.8 of chapter 9:

\medskip
\noindent\textsc{Corollary.}\quad If \(q\in J^+(p)-I^+(p)\), then any causal
curve connecting \(p\) to \(q\) must be a null geodesic.
\medskip

Using similar arguments, we also see that for any set \(S\subseteq M\), we have
\(J^+(S)\subseteq\overline{I^+(S)}\). Since, clearly, \(I^+(S)\subseteq J^+(S)\),
it follows immediately that \(\partial J^+(S)=\partial I^+(S)\). Similarly, we
have \(I^+(S)=\operatorname{int}(J^+(S))\), and hence the boundaries of the
chronological and causal futures of a set are always equal,
\(\partial I^+(S)=\partial J^+(S)\) (see problem 2).

A subset \(S\subseteq M\) is said to be \emph{achronal} if there do not exist
\(p,q\in S\) such that \(q\in I^+(p)\), i.e., if \(I^+(S)\cap S=\emptyset\),
where \(\emptyset\) denotes the empty set. The next theorem asserts that the
boundary of the chronological future of a set always forms a \lq\lq well behaved,\rq\rq\
three-dimensional, achronal surface.

\begin{theorem}\label{thm:8.1.3}
Let \((M,g_{ab})\) be a time orientable spacetime, and let \(S\subseteq M\). Then
\(\partial I^+(S)\) (if nonempty) is an achronal, three-dimensional, embedded,
\(C^0\)-submanifold of \(M\).
\end{theorem}
\begin{proof}
Let \(q\in\partial I^+(S)\). If \(p\in I^+(q)\), then \(q\in I^-(p)\) and since
\(I^-(p)\) is open, an open neighborhood \(O\) of \(q\) is contained in \(I^-(p)\)
as shown in Fig.~\ref{fig:8.3}. Since \(q\) is on the boundary of \(I^+(S)\), we
have \(O\cap I^+(S)\ne\emptyset\) and thus
\(p\in I^+(O\cap I^+(S))\subseteq I^+(S)\). This proves that
\(I^+(q)\subseteq I^+(S)\). Similarly we have \(I^-(q)\subseteq M-I^+(S)\).
Now, if \(\partial I^+(S)\) failed to be achronal, we could find
\(q,r\in\partial I^+(S)\) such that \(r\in I^+(q)\) and hence
\(r\in I^+(S)\). However, this is impossible since \(I^+(S)\) is open and
therefore \(\partial I^+(S)\cap I^+(S)=\emptyset\). This proves that
\(\partial I^+(S)\) is achronal. To obtain the manifold structure of
\(\partial I^+(S)\), we introduce Riemannian normal coordinates \(x^0,x^1,x^2,x^3\)
at \(q\in\partial I^+(S)\) and consider a sufficiently small neighborhood of \(q\)
that \((\partial/\partial x^0)^a\) is everywhere timelike and each of the integral
curves of \((\partial/\partial x^0)^a\) enters \(I^+(q)\subseteq I^+(S)\) and
\(I^-(q)\subseteq M-I^+(S)\). But this implies that each such integral curve
intersects \(\partial I^+(S)\), and since \(\partial I^+(S)\) is achronal, it must
intersect it at precisely one point. Thus,

% Source: Wald GR, physical PDF p. 205
%         (printed p. 192).
% Coverage: A set and the boundary of its chronological future.
% Figures/tables/footnotes: Figure 8.3; no tables or footnotes.
% Status: source-reviewed and compile-clean.
% Uncertainties: none.

\begingroup
\renewcommand{\thefigure}{8.3}
\begin{figure}[H]
  \centering
  \begin{tikzpicture}[x=.88cm,y=.88cm,>=Latex,line cap=round,font=\small]
    \draw[thick] (-3.15,2.00) -- (-1.42,.20)
      .. controls (-.54,.54) and (.34,.26) .. (1.18,.38) -- (2.84,1.88);
    \draw[->] (-1.16,1.18) -- (-2.12,1.02);
    \node at (-.80,1.30) {\(\partial I^+(S)\)};
    \draw[->] (.54,-.10) -- (.43,.35) node[below=.15cm] {\(S\)};
    \fill (1.89,1.35) circle (1.6pt) node[right] {\(q\)};
    \fill (1.91,2.34) circle (1.6pt) node[right] {\(p\)};
    \draw (1.89,1.35) .. controls (1.65,1.11) and (1.54,1.43) .. (1.63,1.58)
      .. controls (1.77,1.79) and (2.10,1.63) .. (2.18,1.39)
      .. controls (2.25,1.19) and (2.08,1.05) .. (1.89,1.35);
    \draw[->] (1.48,1.83) -- (1.67,1.53);
    \node at (1.35,1.95) {\(O\)};
  \end{tikzpicture}
  \caption{A spacetime diagram showing a set \(S\) and the boundary of its future,
  \(\partial I^+(S)\) (see theorem~\ref{thm:8.1.3}).}
  \label{fig:8.3}
\end{figure}
\endgroup


in each such neighborhood, we get a one-to-one association of points of
\(\partial I^+(S)\) with coordinates \((x^1,x^2,x^3)\) characterizing the integral
curve of \((\partial/\partial x^0)^a\). Furthermore, using the achronality of
\(\partial I^+(S)\), the value of \(x^0\) at the intersection point must be a
continuous (in fact, \(C^{1-}\))\footnote{See Hawking and Ellis (1973) for the
definition of \(C^{1-}\).} function of the coordinates \((x^1,x^2,x^3)\), and
thus the above map from a neighborhood of \(q\) in \(\partial I^+(S)\) into
\(\mathbb{R}^3\) is a homeomorphism in the induced topology on \(\partial I^+(S)\).
By repeating this construction for all \(q\in\partial I^+(S)\) we obtain a
\(C^0\)-compatible family of charts covering \(\partial I^+(S)\) which makes
\(\partial I^+(S)\) an embedded submanifold.
\end{proof}

Before proceeding further, we need to introduce several notions that will play an
important role in many considerations of this chapter and of chapter 9. First,
although for many purposes it suffices to consider only differentiable (or even
\(C^\infty\)) curves, when it comes to taking limits (as we shall do below) it is
essential to extend consideration to continuous curves. The definitions of a timelike
or causal curve can be extended to the continuous case by requiring that, locally,
pairs of points on the curve can be joined by a differentiable timelike or,
respectively, causal curve. More precisely, a \emph{continuous curve} \(\lambda\) is
said to be a \emph{future directed timelike (or causal) curve} if for each
\(p\in\lambda\) there exists a convex normal neighborhood \(U\) of \(p\) such that
if \(\lambda(t_1),\lambda(t_2)\in U\) with \(t_1<t_2\), then there exists a future
directed differentiable timelike (or, respectively, causal) curve in \(U\) from
\(\lambda(t_1)\) to \(\lambda(t_2)\). The timelike or causal nature of a continuous
curve clearly is unchanged by a continuous, one-to-one reparameterization, and in
the following we shall consider two curves differing by such a reparameterization to
be equivalent.

Next, we define the notion of extendibility of a continuous curve. We need to
distinguish clearly between the possibilities that a curve \lq\lq runs off to infinity\rq\rq\ or
\lq\lq runs around and around forever\rq\rq\ or \lq\lq runs into a singularity\rq\rq\ as opposed to
the possibility that a curve \lq\lq stops\rq\rq\ somewhere simply because one did not define it
to go further. This distinction can be made precise via the notion of an endpoint of a
curve. Let \(\lambda(t)\) be a future directed causal curve. We say that \(p\in M\)
is a \emph{future endpoint} of \(\lambda\) if for every neighborhood \(O\) of \(p\)
there exists a \(t_0\) such that \(\lambda(t)\in O\) for all \(t>t_0\). (Thus, by
the Hausdorff property of \(M\), \(\lambda\) can have, at most, one future endpoint.
Note also that the endpoint need not lie on the curve, i.e., there need not exist a
value of \(t\) such that \(\lambda(t)=p\).) The curve \(\lambda\) is said to be
\emph{future inextendible} if it has no future endpoint. Past inextendibility is
defined similarly. Note that if \(\lambda\) is a differentiable causal curve with future
endpoint \(p\), then it may not be possible to extend \(\lambda\) beyond \(p\) as a
differentiable causal curve, but \(\lambda\) always can be extended as a continuous
causal curve by adjoining a continuous causal curve to \(\lambda\) at \(p\).

An important technical lemma relating to extendibility of causal curves which will
be used in section 8.3 is the following.

\begin{lemma}\label{lem:8.1.4}
Let \(\lambda\) be a past inextendible causal curve passing through point \(p\).
Then through any \(q\in I^+(p)\) there exists a past inextendible timelike curve
\(\gamma\) such that \(\gamma\subseteq I^+(\lambda)\).
\end{lemma}
\begin{proof}
Without loss of generality we may assume that the curve parameter, \(t\), of
\(\lambda\) has the range \([0,\infty)\). Following Geroch (1970b), we choose an
arbitrary Riemannian metric on \(M\). Using theorem~\ref{thm:8.1.2} we can
construct a timelike curve \(\gamma(t)\) for \(t\in[0,1]\) which starts at \(q\) and
satisfies the property that \(\gamma\subseteq I^+(\lambda)\) and
\(d(\gamma(t),\lambda(t))<C/(1+t)\), where \(C\) is a constant and the distance,
\(d\), is the greatest lower bound of the length, measured using the Riemannian
metric, of all curves connecting a point \(\gamma(t)\) to a point \(\lambda(t')\) with
\(t,t'\in[0,1]\). By induction, we continue to extend \(\gamma(t)\) to a curve
defined for \(t\in[0,n]\) for all \(n\) (i.e., for \(t\in[0,\infty)\)) preserving
these properties. The resulting \(\gamma\) will be past inextendible because any
endpoint of \(\gamma\) would also be an endpoint of \(\lambda\). Thus, \(\gamma\) is
the desired curve.
\end{proof}

Another notion that will play a prominent role in many arguments of this chapter
and of chapter 9 is the convergence of causal curves. Let \(\{\lambda_n\}\) be a
sequence of causal curves. A point \(p\in M\) is said to be a \emph{convergence
point} of \(\{\lambda_n\}\) if, given any open neighborhood \(O\) of \(p\), there
exists an \(N\) such that \(\lambda_n\cap O\ne\emptyset\) for all \(n>N\). A curve
\(\lambda\) is said to be a \emph{convergence curve} of \(\{\lambda_n\}\) if each
\(p\in\lambda\) is a convergence point. Similarly, \(p\) is said to be a \emph{limit
point} of \(\{\lambda_n\}\) if every open neighborhood of \(p\) intersects infinitely
many \(\lambda_n\). A curve \(\lambda\) is said to be a \emph{limit curve} of
\(\{\lambda_n\}\) if there exists a subsequence \(\{\lambda'_n\}\) for which
\(\lambda\) is a convergence curve. (Thus, if \(\lambda\) is a limit curve, then each
\(p\in\lambda\) is a limit point. Note, however, that a curve \(\gamma\) such that
each \(p\in\gamma\) is a limit point need not be a limit curve.) The following result
plays a crucial role in many of the proofs given below.

\begin{lemma}\label{lem:8.1.5}
Let \(\{\lambda_n\}\) be a sequence of future inextendible causal curves which have
a limit point \(p\). Then there exists a future inextendible causal curve \(\lambda\)
passing through \(p\) which is a limit curve of the \(\{\lambda_n\}\).
\end{lemma}
\begin{proof}[Sketch of proof]
We choose a convex normal neighborhood of \(p\) and a ball of Riemannian normal
coordinate radius \(R\) about \(p\) contained in this neighborhood. We pick a
subsequence of \(\{\lambda_n\}\) that converges to \(p\) and, using the compactness
of the coordinate sphere of radius \(R\), a sub-subsequence which converges to a
point on this sphere. Then we examine, in turn, all the coordinate spheres whose
radii are rational multiples, between 0 and 1, of \(R\) and continue to extract limit
points lying on these spheres and subsequences converging to these points. Finally,
we take the closure of the set of these limit points and show that it defines a
continuous causal limit curve \(\lambda\). We then go to the endpoint of \(\lambda\)
on the sphere of radius \(R\) and repeat the procedure. Continuing in this manner,
we extend \(\lambda\) indefinitely.\footnote{It should be noted that the above
argument shows the existence of an extendible limit curve through \(p\), together
with the fact that any extendible limit curve can be further extended as a limit curve.
To assert existence of an inextendible limit curve (i.e., a maximal element in the
set of limit curves under ordering by inclusion) by this argument we must appeal to
Zorn's lemma. However, it appears that a proof of lemma 8.1.5 could be given
without invoking Zorn's lemma.} Technical details of this proof can be found in
Hawking and Ellis (1973).
\end{proof}

It should be noted that in theorem 8.1.5 if each \(\lambda_n\) is a timelike curve,
the limit curve \(\lambda\) still may be only a causal curve, i.e., a sequence of
timelike curves may converge to a null curve. Similarly, even if all the
\(\lambda_n\) are smooth curves, \(\lambda\) may be only continuous.

As a direct application of theorem 8.1.5, we prove, now, a theorem characterizing
the nature of boundaries of chronological futures.

\begin{theorem}\label{thm:8.1.6}
Let \(C\) be a closed subset of the spacetime manifold \(M\). Then every point
\(p\in\partial I^+(C)\) with \(p\notin C\) lies on a null geodesic \(\lambda\)
which lies entirely in \(\partial I^+(C)\) and either is past inextendible or has a
past endpoint on \(C\).
\end{theorem}
\begin{proof}
Choose a sequence of points \(\{q_n\}\) in \(I^+(C)\) which converges to \(p\), as
illustrated in Fig.~\ref{fig:8.4}. For each \(q_n\) let \(\lambda_n\) be a past
directed timelike curve connecting \(q_n\) to a point in \(C\). Consider, now, the
new spacetime manifold \(M-C\) obtained by removing the set \(C\). (It is here that
we use the assumption that \(C\) is closed, for otherwise \(M-C\) would not define a
manifold.) On \(M-C\), each \(\lambda_n\) is past inextendible and \(p\) is a limit
point of the sequence \(\{\lambda_n\}\). Hence, by lemma~\ref{lem:8.1.5} there
exists a past inextendible causal limit curve \(\lambda\) passing through \(p\). Each
point of \(\lambda\) is a limit point of sequences in \(I^+(C)\), so
\(\lambda\subseteq\overline{I^+(C)}\). On the other hand, if any point of
\(\lambda\) were in \(I^+(C)\), then by the corollary to theorem~\ref{thm:8.1.2}
we would have \(p\in I^+(C)\), since \(p\) could be connected to \(C\) by a causal
curve which is not a null geodesic. This contradicts the fact that
\(p\in\partial I^+(C)\). Thus, \(\lambda\subseteq\partial I^+(C)\). Furthermore,
the achronality of \(\partial I^+(C)\) (theorem~\ref{thm:8.1.3}), together with the
corollary to theorem~\ref{thm:8.1.2}, implies that \(\lambda\) must be a null
geodesic. Finally, since \(\lambda\) is past inextendible in \(M-C\), in \(M\) it
must either remain past inextendible or have a past endpoint on \(C\).
\end{proof}

% Source: Wald GR, physical PDF p. 208
%         (printed p. 195).
% Coverage: A sequence in the chronological future of a closed set converging to its boundary.
% Figures/tables/footnotes: Figure 8.4; no tables or footnotes.
% Status: source-reviewed and compile-clean.
% Uncertainties: none.

\begingroup
\renewcommand{\thefigure}{8.4}
\begin{figure}[H]
  \centering
  \begin{tikzpicture}[x=.82cm,y=.82cm,>=Latex,line cap=round,font=\small]
    \draw[thick] (-2.55,.10) .. controls (-1.55,.58) and (-.42,.10) .. (.90,.20)
      -- (2.72,2.42);
    \node at (1.16,1.60) {\(I^+(C)\)};
    \draw[->] (.22,-.08) -- (.14,.32) node[below=.15cm] {\(C\)};
    \fill (-1.98,2.30) circle (1.6pt) node[left] {\(p\)};
    \foreach \x/\y in {-1.72/2.64,-1.57/2.98,-1.42/3.33,-1.30/3.65} {
      \fill (\x,\y) circle (1.25pt);
    }
    \node at (-1.55,3.92) {\(q_n\)};
    \draw[thick] (-2.55,.10) .. controls (-2.65,1.32) and (-2.33,2.26) .. (-1.72,2.64);
    \draw[thick] (-2.39,.14) .. controls (-2.34,1.50) and (-2.08,2.42) .. (-1.57,2.98);
    \draw[thick] (-2.22,.16) .. controls (-2.08,1.65) and (-1.80,2.70) .. (-1.42,3.33);
    \draw[thick] (-2.03,.18) .. controls (-1.82,1.75) and (-1.59,2.93) .. (-1.30,3.65);
    \draw[->] (-1.10,2.42) -- (-1.78,2.12);
    \node at (-.67,2.54) {\(\lambda_n\)};
  \end{tikzpicture}
  \caption{A spacetime diagram showing a sequence of points in \(I^+(C)\)
  converging to \(p\in\partial I^+(C)\) (see theorem~\ref{thm:8.1.6}).}
  \label{fig:8.4}
\end{figure}
\endgroup


An example where \(\lambda\) is past inextendible is provided by point \(q\) in
Fig.~\ref{fig:8.2}. Note that although \(\lambda\) cannot have a past endpoint
except on \(C\), it may have a future endpoint; an example where \(\lambda\) has a
future endpoint is given in Fig.~\ref{fig:8.5}.

% Source: Wald GR, physical PDF p. 208
%         (printed p. 195).
% Coverage: A flat two-dimensional spacetime with topology \(\mathbb{R}\times S^1\).
% Figures/tables/footnotes: Figure 8.5; no tables or footnotes.
% Status: source-reviewed and compile-clean.
% Uncertainties: none.

\begingroup
\renewcommand{\thefigure}{8.5}
\begin{figure}[H]
  \centering
  \begin{tikzpicture}[x=.9cm,y=.9cm,>=Latex,line cap=round,font=\small]
    \draw[thick] (-1.55,3.65) arc[start angle=180,end angle=360,x radius=1.55,y radius=.35];
    \draw[thick] (-1.55,3.65) arc[start angle=180,end angle=0,x radius=1.55,y radius=.35];
    \draw[thick] (-1.55,.30) arc[start angle=180,end angle=360,x radius=1.55,y radius=.35];
    \draw[dashed] (-1.55,.30) arc[start angle=180,end angle=0,x radius=1.55,y radius=.35];
    \draw[thick] (-1.55,.30) -- (-1.55,3.65);
    \draw[thick] (1.55,.30) -- (1.55,3.65);
    \fill (0,.88) circle (1.6pt) node[below] {\(p\)};
    \fill (0,2.82) circle (1.6pt) node[above] {\(q\)};
    \draw[thick] (0,.88) .. controls (-.62,1.34) and (-1.46,1.68) .. (-1.55,2.03);
    \draw[dashed] (-1.55,2.03) .. controls (-.82,2.45) and (-.30,2.66) .. (0,2.82);
    \draw[thick] (0,.88) .. controls (.63,1.34) and (1.46,1.68) .. (1.55,2.03);
    \draw[dashed] (1.55,2.03) .. controls (.82,2.45) and (.30,2.66) .. (0,2.82);
  \end{tikzpicture}
  \caption{A spacetime diagram of a two-dimensional flat spacetime with topology
  \(\mathbb{R}\times S^1\). The null geodesic generators of \(\partial I^+(p)\)
  have \(q\) as a future endpoint.}
  \label{fig:8.5}
\end{figure}
\endgroup

